\clearpage\thispagestyle{otFront}\phantomsection\label{mode-emploi}\pdfbookmark[0]{Mode d’emploi}{bm-mode-emploi}\otPageTitre{Mode d’emploi}{Des repères pour travailler seul ou accompagné}\fontsize{11}{15}\selectfont

\begin{tabularx}{\linewidth}{@{}lY@{}}\toprule
Une pastille & Découverte : lire, écrire, calculer directement.\\[3mm]
Deux pastilles & Entraînement : choisir une méthode, organiser ses étapes.\\[3mm]
Trois pastilles & Problèmes : raconter, modéliser, calculer et répondre.\\\bottomrule
\end{tabularx}\par\medskip
Ces pastilles décrivent les tâches ; elles ne classent pas les élèves.\par\medskip
\begin{exemple}[Lire une correction]Pour « $47+\square=92$ », on cherche le
nombre qui complète 47 jusqu’à 92 : c’est $92-47=45$. Le contrôle
$47+45=92$ confirme. Refaire ensuite le calcul sans regarder
aide plus que recopier la réponse.\end{exemple}
\begin{attention}Lire l’énoncé en entier avant de calculer. Un mot comme
« plus » ou « de moins » ne décide pas de l’opération : c’est l’histoire
du problème qui la choisit. Toujours écrire l’unité dans la réponse.\end{attention}
Chercher d’abord, ouvrir le corrigé ensuite. Après lecture, refaire le
passage difficile sur une feuille : c’est là que l’on apprend.

\clearpage
